\documentclass[11pt]{article}
\usepackage[T1]{fontenc}
\usepackage[utf8]{inputenc}
\usepackage{graphicx,url,xspace,amsmath,amsfonts,booktabs,enumitem}
\usepackage[hidelinks]{hyperref}
\setlength{\oddsidemargin}{.25in}
\setlength{\evensidemargin}{.25in}
\setlength{\textwidth}{6in}
\setlength{\topmargin}{-0.4in}
\setlength{\textheight}{8.5in}
\setlist{nosep}

\def\subjnum{601.445/601.645}
\def\subjname{Practical Cryptographic Systems}

%
% Impressive box to go at the top of the first page
%
\newcommand{\handout}[5]{
  \renewcommand{\thepage}{#1-\arabic{page}}
  \noindent
  \begin{center}
    \framebox{
      \vbox{
        \hbox to 5.78in { {\bf \subjnum: \subjname} \hfill #2 }
        \vspace{4mm}
        \hbox to 5.78in { {\Large \hfill #5  \hfill} }
        \vspace{2mm}
        \hbox to 5.78in { {\it #3 \hfill #4} }
        }
      }
  \end{center}
  \vspace*{4mm}
  }

\newcommand{\code}[1]{\texttt{#1}}
\newcommand{\file}[1]{\texttt{#1}}

\emergencystretch=1.5em
\begin{document}
\handout{Anon Report}{September 2, 2026}{Instructor: Matthew Green}{Applies to Assignments 1, 2, and 3}{Submitting Your Anonymous Agent Report}

\noindent
This handout explains how to turn in the \emph{agent report}. This is the short, anonymous
write-up your coding agent produces about how well its teaching instructions worked. The tool,
\file{tools/anon-report.py}, is one file of plain Python with no dependencies; you can
run every step yourself, or let your agent do it (\file{AGENTS.md} contains these instructions as well.)

\section*{What it is, and why it works this way}

The report is feedback on the \emph{instructions} not on your performance. It tells us what the agent did
versus what you did, where the ``explain first'' rule helped or got in the way, and
what it would change. This file is not graded. We want it to be candid, so it's anonymous ---
and rather than ask you to trust our promise on that, the system uses cryptography so that we can't
link a report to a student. If you're interested, here's what it offers:

\medskip \noindent
\begin{itemize}
\item \textbf{Encryption.} Your report is encrypted on your own machine to the course
  public key. The drop box and the public bulletin board hold only ciphertext, and only the
  instructor's offline key can decrypt it.
\item \textbf{A ring signature.} Each report is signed with a \emph{linkable ring
  signature} over the roster of registered student keys (your Part 0 key). It proves
  ``an enrolled student wrote this'' without saying which one, and it carries a tag
  that makes a second deposit from the same key visible as a duplicate, but without
  revealing the key. Tags differ per assignment, so your reports across the semester
  are not linkable to each other either.
\item \textbf{Tor.} The deposit travels over Tor to an onion service, so the box
  never sees your IP address. The application keeps no per-request log and zeroes the
  stored file's modification time; ordinary system metadata can still reflect when it arrived.
\end{itemize}

\medskip \noindent
\noindent Later in the course you might implement this scheme yourself; it is the
same math as the public-key block, used for a privacy-preserving and practical purpose.

\section*{What it doesn't protect}

Cryptography can offer only so much security. Filesystem or network records may
reflect when you connect, and a distinctive writing style could narrow things down.
We do not use that metadata to identify authors, but we will read the reports.
The agent is instructed to write the report
without names, JHED, email, file paths, verbatim code, or other identifying details,
and to show it to you first. \textbf{Read the report before it is sent.} You're the
last check: you can edit it, although we'd prefer you not remove any interesting content.

\section*{Dates (Assignment 1)}

\begin{center}\small
\begin{tabular}{ll}
\toprule
Fri Sep 11 & Part 0: register your public key (Gradescope ``A1 Part 0'') \\
Mon Sep 14 & Roster published: \file{anon-report/roster.json} appears in the course repository \\
by Sun Sep 20 & Deposit your report \\
$\approx$ Mon Sep 21 & Tags snapshot published: \file{anon-report/tags-A1.json} \\
by Fri Sep 25 & Attendance proof (Gradescope ``A1 attendance'') \\
\bottomrule
\end{tabular}
\end{center}
For Assignment 1 the channel is a \textbf{pilot}: the attendance proof is not required
for review-lab sign-up. From Assignment 2 on, it will be.

\section*{Step by step}

You can have your agent submit the report. Honestly, I would --- the process is a bit annoying. If you choose to do it yourself, all of the 
instructions are below.  
Every command below is run from your assignment directory, on the machine that holds
your Part 0 key (\file{\textasciitilde/.anon-report/key.json}).

\medskip \noindent
\begin{enumerate}
\item \textbf{Get the roster.} On or about \textbf{Mon Sep 14} the course staff publish
  \file{anon-report/roster.json} in the course repository
  (\url{https://github.com/matthewdgreen/practicalcrypto2026}). It contains every
  registered key, the course encryption key, and the drop box's onion address. Pull the
  repository (or download that one file) so your local copy is current. If the file isn't
  there yet, the channel isn't open.
\item \textbf{Before you deposit.} Have \textbf{Tor Browser} open and connected (free,
  \url{https://www.torproject.org}; the tool talks to it on port 9150 --- a running
  \code{tor} daemon on 9050 also works).
\item \textbf{Read your report.} Have it as a text file, e.g.\ \file{agent-report.md}.
  Check it contains no name, JHED, email, paths, or code. This is the last moment it's
  editable.
\item \textbf{Deposit it} --- encrypts, signs, and uploads over Tor in one step:
\begin{verbatim}
  python3 tools/anon-report.py deposit agent-report.md \
      --roster anon-report/roster.json \
      --assignment A1
\end{verbatim}
  Success looks like \code{ACCEPTED   tag 3F9A...}. Note the tag: it's the name of
  your entry on the public bulletin board. A receipt is saved under
  \file{\textasciitilde/.anon-report/receipts/}.
\item \textbf{If the upload fails} (typically ``Tor not connected''), don't re-run
  \code{deposit}. Fix Tor --- on campus wifi, Tor Browser $\to$ Settings $\to$ Connection
  $\to$ Bridges $\to$ Snowflake --- then re-send the saved receipt:
\begin{verbatim}
  python3 tools/anon-report.py upload \
      ~/.anon-report/receipts/<newest file> \
      --roster anon-report/roster.json
\end{verbatim}
\item \textbf{One shot.} The drop box keeps the \emph{first} deposit per key per
  assignment and refuses a second as a duplicate; there are no amendments. If staff
  publish a corrected roster, update it and deposit again: the server can replace only
  your same-tag entry from the superseded roster. Never use \code{-{}-direct}: it skips
  Tor and shows the server your IP address.
\item \textbf{Optional: check your entry} once the board is published:
\begin{verbatim}
  python3 tools/anon-report.py verify-bundle \
      anon-report/bulletin/A1/<your tag>.json \
      --roster anon-report/roster.json
\end{verbatim}
\item \textbf{Attendance} (after \file{anon-report/tags-A1.json} appears). This proves
  you deposited \emph{a} report without revealing which; your nonce is your JHED:
\begin{verbatim}
  python3 tools/anon-report.py attend \
      --roster anon-report/roster.json --assignment A1 \
      --tags anon-report/tags-A1.json --nonce <your JHED>
\end{verbatim}
  Upload the resulting \file{attendance.json} to the Gradescope assignment
  \textbf{``A1 attendance''}, under your own name. If the tool says your tag isn't in
  the list, your deposit didn't land or the snapshot predates it --- ask on Piazza
  rather than using \code{-{}-force}.
\item \textbf{Keep things separate.} The report, the receipt, and
  \file{attendance.json} never go into your main assignment submission on Gradescope.
  Keep \file{\textasciitilde/.anon-report/} for the rest of the semester.
\end{enumerate}

\section*{Questions you might have}

\begin{itemize}
\item \emph{I didn't use an agent.} Then there is no report to deposit for this
  assignment; nothing else to do.
\item \emph{I lost my key.} You can't sign or prove attendance without it. Generate a
  new one with \code{keygen} and tell us on Piazza (privately); it goes on the next
  roster version. Back the key up this time.
\item \emph{Tor won't connect on campus.} Use a bridge: Tor Browser $\to$ Settings
  $\to$ Connection $\to$ Bridges $\to$ Snowflake. Nothing else changes.
\item \emph{Can I see that my report arrived?} Yes --- step 7. The board is public
  precisely so that everyone can check the box isn't dropping deposits.
\item \emph{Can the instructor tell it's mine?} Not from the system. See ``What it
  doesn't protect'' for the honest residual: timing and writing style.
\item \emph{What does the server keep?} Ciphertext bundles, their ring signatures, and
  tags. No client IP addresses, plaintext, or per-request application log; ordinary
  system metadata may still reflect receipt timing.
\end{itemize}


\bigskip
\noindent\rule{\textwidth}{0.4pt}

\medskip\noindent\textbf{Authorship and AI usage note} \emph{(in the same format this course asks of you).} This handout and the initial tool and server were written by \textbf{Matthew Green} with \textbf{Claude} (Anthropic's Fable~5 model, via Claude Code). \textbf{OpenAI Codex} (the Daybreak Blue and GPT~5.6 Sol models) later performed a security review, hardening, adversarial and Tor end-to-end tests, and deployment verification. An early version used a wrong 2048-bit group prime; the current tool embeds the published RFC~3526 value and checks basic group invariants before use. Verify your constants.

\end{document}
